\documentclass[../../../include/open-logic-section]{subfiles}

\begin{document}

\olfileid{sth}{choice}{countablechoice}
\olsection{انتخابِ شمارا}

بدون هیچ استفاده‌ای از انتخاب یا خوش‌ترتیبی، به‌آسانی می‌توان اثبات
کرد که:

\begin{lem}[در $\Zminus$]
هر مجموعهٔ متناهی دارای تابع انتخاب است.
\end{lem}

\begin{proof}
بگذارید $a = \{b_1, \ldots, b_n\}$. برای سادگی فرض کنید هر
$b_i \neq \emptyset$. پس اشیایی مانند
$c_1, \ldots, c_n$ وجود دارند به‌طوری‌که
$c_1 \in b_1, \ldots, c_n \in b_n$. اکنون بنا بر
\olref[z][pairs]{prop:pairsconsequences}، مجموعهٔ
$\{\langle b_1,c_1\rangle, \ldots, \langle b_n,c_n\rangle\}$ وجود
دارد؛ و این تابعی انتخاب برای~$a$ است.
\end{proof}

اما به‌محض آن‌که مجموعه‌های نامتناهی را در نظر بگیریم، اوضاع مبهم‌تر
می‌شود. برای نمونه، این گسترش «کمینه» از نتیجهٔ بالا را در نظر بگیرید:

\begin{defish}
\emph{انتخاب شمارا.} هر مجموعهٔ \emph{شمارا} دارای تابع انتخاب است.
\end{defish}

این حالت خاصی از انتخاب است. روشن می‌شود که این اصل نسبتاً بسیار،
بی‌آن‌که استفاده از آن آشکارا تشخیص داده شود، به کار رفته است. دو
نمونهٔ خوب را ببینید.\footnote{این نمونه‌ها از
\citet[\S9.4]{Potter2004} و لوکا اینکورواتی‌اند.}

\begin{ex}
اندیشه‌ای طبیعی چنین است: برای هر مجموعهٔ $A$، یا
$\cardle{\omega}{A}$، یا برای یک $n \in \omega$ داریم
$\cardeq{A}{n}$. این یکی از راه‌های بیان این تصور شهودی است که هر
مجموعه یا متناهی است یا نامتناهی. کانتور و بسیاری از ریاضی‌دانان
دیگر این ادعا را بی‌اثبات مطرح کردند. ما که محتاطیم، آن را در
\olref[cardinals][classing]{generalinfinitycharacter} اثبات کردیم.
اما در آن اثبات در $\ZFC$ کار می‌کردیم، زیرا فرض کرده بودیم هر
مجموعهٔ $A$ را می‌توان خوش‌ترتیب کرد و در نتیجه وجودِ $\card{A}$
تضمین می‌شود. یعنی انتخاب را به‌صراحت فرض کرده بودیم.

در واقع، \citet{Dedekind1888} اثبات خود را از این ادعا چنین عرضه
کرد:

\begin{thm}[در $\Zminus + \text{انتخاب شمارا}$]
برای هر $A$، یا $\cardle{\omega}{A}$، یا برای یک
$n \in \omega$ داریم $\cardeq{A}{n}$.
\end{thm}

\begin{proof}
فرض کنید برای همهٔ $n \in \omega$ داشته باشیم
$\cardneq{A}{n}$. در این صورت، به‌ویژه برای هر $n < \omega$
زیرمجموعه‌ای مانند $A_n \subseteq A$ با دقیقاً
$\cardexpo{2}{n}$ عضو وجود دارد. با استفاده از دنبالهٔ
$A_0, A_1, A_2, \ldots$، برای هر $n$ تعریف می‌کنیم:
\[
	B_n = A_n \setminus \bigcup_{i < n} A_i.
\]
اکنون توجه کنید:
\begin{align*}
	\card{\bigcup_{i < n}A_i}
	&\leq \card{A_0} + \card{A_1} + \ldots + \card{A_{n-1}}\\
	&=1 + 2 + \ldots + 2^{n-1}\\
	& = 2^n - 1\\
	& < 2^n = \card{A_n}
\end{align*}
پس هر $B_n$ دست‌کم یک عضو، مانند $c_n$، دارد. افزون بر این،
$B_n$ها دوبه‌دو مجزایند؛ پس اگر $c_n = c_m$، آن‌گاه $n = m$.
اما هر $c_n \in A$. بنابراین تابع $f(n) = c_n$ تزریقی
$\omega \to A$ است.
\end{proof}
\noindent
ددکیند یادآوری نکرد که انتخاب شمارا را به کار برده است. اما
\emph{شما} استفاده از آن را دیدید؟ دوباره نگاه کنید. (واقعاً:
\emph{دوباره نگاه کنید}.)

اثبات دو بار از انتخاب شمارا استفاده کرد. بار نخست آن را برای
به‌دست‌آوردن دنبالهٔ مجموعه‌های
$A_0$، $A_1$، $A_2$، \dots\@ به کار بردیم. سپس بار دیگر برای
برگزیدن اعضای $c_n$ از هر~$B_n$ از آن استفاده کردیم. افزون بر این،
این استفاده از انتخاب حذف‌ناپذیر است. \citet[p.~138]{Cohen1966}
اثبات کرد که بدون هیچ صورتی از انتخاب، نتیجه شکست می‌خورد. یعنی
وجود مجموعه‌هایی \emph{مقایسه‌ناپذیر} با~$\omega$ با $\ZF$ سازگار
است.
\end{ex}

\begin{ex}
کانتور در \citeyear{Cantor1878} اظهار کرد که اجتماع شمارایی از
مجموعه‌های شمارا، شماراست. او اثباتی عرضه نکرد؛ شاید چون آن را بدیهی
می‌دانست. ما که محتاطیم، صورت کلی‌تری از این نتیجه را در
\olref[card-arithmetic][simp]{kappaunionkappasize} اثبات کردیم. اما
اثبات ما به‌صراحت انتخاب را فرض می‌کرد. حتی اثبات نتیجهٔ کم‌کلی‌تر
نیز به انتخاب شمارا نیاز دارد.

\begin{thm}[در $\Zminus + \text{انتخاب شمارا}$]
اگر برای هر $n \in \omega$، مجموعهٔ $A_n$ شمارا باشد، آن‌گاه
$\bigcup_{n < \omega} A_n$ شماراست.
\end{thm}

\begin{proof}
بدون کاستن از کلیت، فرض کنید هر $A_n \neq \emptyset$. پس برای هر
$n \in \omega$، !!a{surjection} مانند
$f_n \colon \omega \to A_n$ وجود دارد. تابع
$f \colon \omega \times \omega \to \bigcup_{n < \omega} A_n$ را با
$f(m, n) = f_n(m)$ تعریف کنید. نتیجه از شمارابودن
$\omega \times \omega$
(\olref[sfr][siz][zigzag]{natsquaredenumerable}) و !!a{surjection}
بودنِ $f$ به دست می‌آید.
\end{proof}
\noindent
آیا کاربرد انتخاب شمارا را دیدید؟ از آن برای برگزیدن دنبالهٔ توابع
$f_0$، $f_1$، $f_2$، \dots استفاده کردیم.\footnote{کاربرد مشابهی
از انتخاب در \olref[card-arithmetic][simp]{kappaunionkappasize} رخ
داد، آنجا که دستور دادیم: «برای هر
$\beta \in \cardfont{a}$، یک !!a{injection} مانند
$f_\beta$ را ثابت کنید.»} باز هم، در نبود هرگونه اصل انتخاب، نتیجه
شکست می‌خورد. به‌طور مشخص، \citet{FefermanLevy1963} اثبات کرد که
این گزاره با $\ZF$ سازگار است: اجتماع شمارایی از مجموعه‌های شمارا
دارای کاردینالیتهٔ~$\beth_1$ باشد. اما راسل این نکته را بسیار
بامزه‌تر بیان می‌کند:
\begin{quote}
  این را می‌توان با میلیونری نشان داد که هرگاه یک جفت چکمه می‌خرید،
  یک جفت جوراب نیز می‌خرید و هرگز در زمان دیگری چیزی از این دو
  نمی‌خرید؛ و چنان شیفتهٔ خرید هر دو بود که سرانجام
  $\aleph_0$ جفت چکمه و $\aleph_0$ جفت جوراب داشت\dots\@ در میان
  چکمه‌ها می‌توانیم راست و چپ را از هم بازشناسیم و بنابراین از هر
  جفت یکی را انتخاب کنیم: همهٔ چکمه‌های راست یا همهٔ چکمه‌های چپ را.
  اما دربارهٔ جوراب‌ها هیچ اصل انتخاب مشابهی خود را پیشنهاد نمی‌کند،
  و جز با فرض اصل ضربی [یعنی در عمل اصل انتخاب] نمی‌توانیم مطمئن
  باشیم کلاسی هست که از هر جفت یک جوراب در آن باشد.
  \citep[p.~126]{Russell1919}
\end{quote}
خلاصه اینکه برای اثبات گزارهٔ زیر به صورتی از انتخاب نیاز داریم:
اگر شمارا جفت جوراب داشته باشید، آن‌گاه (فقط) شمارا جوراب دارید.
در واقع، بدون انتخاب شمارا (یا اصلی هم‌ارز با آن)، ممکن است اجتماع
شمارایی از مجموعه‌های شمارا، شمارا نباشد.
\end{ex}

درس ماجرا این است که انتخاب شمارا بارها به کار رفته، بی‌آن‌که
کاربرانش چندان از آن آگاه باشند. پرسش فلسفی این است: چگونه می‌توانیم
انتخاب شمارا را \emph{توجیه} کنیم؟

یک کوشش برای توجیه شهودی شاید به اَبَروظیفه‌ای متوسل شود. فرض کنید
انتخاب نخست را در $\nicefrac{1}{2}$ دقیقه، انتخاب دوم را در
$\nicefrac{1}{4}$ دقیقه، \dots، انتخاب $n$ام را در
$\nicefrac{1}{2^n}$ دقیقه، \dots انجام دهیم. آن‌گاه در مدت
$1$~دقیقه، دنباله‌ای $\omega$تایی از انتخاب‌ها را انجام داده و یک
تابع انتخاب تعریف کرده‌ایم.

اما چنین آزمایش فکری واقعاً چه چیزی می‌تواند به ما بیاموزد؟ نخست
اینکه بر برداشتی بیش‌ازحد لفظی از «انتخاب‌کردن» تکیه دارد. دیگر
اینکه به نظر می‌رسد ریاضیات را به امکان متافیزیکی گره می‌زند.

مهم‌تر آن‌که این اندیشه هیچ توجیهی برای انتخاب \emph{به‌تمام‌معنا}،
در برابر \emph{صرفِ} انتخاب شمارا، فراهم نمی‌کند. زیرا اگر بخواهیم
\emph{هر} مجموعه تابع انتخاب داشته باشد، باید بتوانیم
«اَبَروظیفه‌ای با هر طول ترتیبی» انجام دهیم. بی‌پرده بگوییم، این
اندیشه خنده‌آور است.

\end{document}
